\section{Related work}
The analytical study of the NTK at infinite width has let to numerous insights about the training dynamics with results on trainability~\cite{xiao2020}, double-descent~\cite{geiger2020a}, generative adversarial networks~\cite{franceschi2022}, physics-informed neural networks~\cite{wang2022b}, data augmentation~\cite{gerken2024} and equivariant neural networks~\cite{misof2025}, among others. Infinite-width NTKs have also been used to optimize various aspects of the training process, such as architecture selection and bias detection~\cite{deoliveira2025}.

In the machine-learning literature, finite-width corrections in the formulation used in this article were studied in~\cite{yaida2020}. A book on the subject which develops the background step-by-step and contains many novel results is also available~\cite{roberts2022}. The analogy between QFT and neural networks has been explored from different perspectives in the physics literature. It was studied systematically for the first time in~\cite{halverson2021}. This work also established a connection to renormalization group flow which was studied in a non-perturbative setting in \cite{erbin2022}. In~\cite{grosvenor2022}, a connection to the $O(N)$ vector model was established. Different perturbations of the infinite-width limit related to the Edgeworth expansion have been studied as well~\cite{demirtas2024}. This line of work has led to results about symmetries~\cite{maiti2021}, initialization stability~\cite{banta2024}, orthogonal initializations~\cite{day2023} and scaling laws~\cite{maloney2022, zhang2025} among others.

Feynman diagrams appear in~\cite{banta2024, halverson2021, grosvenor2022, demirtas2024, maloney2022, zhang2025}, but are restricted to preactivation-statistics. In contrast, we use Feynman diagrams to compute NTK and joint prectivation-NTK statistics. This constitutes a substantial extension since preactivation statistics can only describe the initialization of the network, while NTK statistics are needed for training. Also \cite{dyer2019} uses Feynman diagrams to study corrections to the NTK, but their treatment is restricted to providing the scaling behavior of correlation functions, whereas our framework can evaluate the correlation functions explicitly.

%%% Local Variables:
%%% mode: LaTeX
%%% TeX-master: "ntk_feynman"
%%% End:
